\documentclass[11pt]{article}
\usepackage[margin=0.8in]{geometry}
\usepackage[T1]{fontenc}
\usepackage[utf8]{inputenc}
\usepackage{newtxtext,newtxmath}
\usepackage{amsmath,booktabs,array,longtable,graphicx,microtype,xurl}
\usepackage[colorlinks=true,allcolors=blue]{hyperref}
\usepackage{caption}
\makeatletter
\let\tableinput\@@input
\makeatother
\captionsetup{font=small,labelfont=bf}
\renewcommand{\thetable}{S\arabic{table}}
\renewcommand{\thepage}{S\arabic{page}}
\setlength{\parindent}{0pt}
\setlength{\parskip}{0.55em}
\emergencystretch=3em
\newcommand{\metrictable}[2]{%
\begin{table}[htbp]\centering
\caption{#1}
\begin{tabular}{lrrrr}\toprule
Cohort & AUROC & ECE & Brier & Log-loss \\
\midrule\tableinput{tables/#2}\bottomrule
\end{tabular}\end{table}}
\begin{document}
\begin{center}
{\Large\bfseries Supporting Information}\par
\vspace{0.5em}
{\large Random cross-validation understates cross-cohort calibration error in
multi-cohort proteomics: an accession-aligned evaluation-protocol contrast}\par
\vspace{0.5em}
Zeyu Fu\par
\href{mailto:fuzeyu99@126.com}{fuzeyu99@126.com}
\end{center}

\section*{S1. Scope and corrected identities}

This supplement accompanies the accession-aligned six-cohort analysis.
Its tables derive from the same stored predictions used by the article
figures. Historical submission versions and historical gene-indexed
analyses are not updated copies of this supplement.

The 1,191 samples comprise 689 tumors, 496 adjacent-normal samples and
six HNSCC cored-normal samples. The latter are identified by the source
suffix \texttt{.C} and, like \texttt{.N} samples, have label zero.
Both suffixes map to the corresponding unsuffixed patient identifier,
with cohort included in the patient key. The resulting 719 patients
include 470 with both tissue classes, 219 tumor-only and 30 normal-only
patients. HNSCC contributes 110 tumors and 72 normals from 112 patients.
Normal here means a source-defined nonmalignant sample, not a healthy
donor. Patients with tumor, adjacent-normal and cored-normal samples
stay together under the grouped split.

The sample manifest records cohort, sample ID, tissue class, binary
label, true patient ID, and fold for each of the three protocols.
Sample CV splits 381 patients across evaluation folds. Patient-grouped
CV and LOGO split zero true patients. The corrected GroupKFold
partition was generated in scikit-learn 1.9.0 and is preserved by sample
identity; it is not the historical partition based on 723 parsed keys.

The primary protein intersection comprises 8,906 stable ENSP accessions.
Version suffixes are stripped only after verifying unique stable
accessions per cohort; gene annotation and, for the retained set, the
complete accession version also agree across all cohorts. There is no
gene-level averaging in this primary matrix. An old 9,675-gene
intersection selected the first protein for each name; 1,234 names
selected different stable accessions in different cohorts. The
8,441-name sensitivity retains only the consistent first-accession
subset. Neither mapping guarantees absence of every protein-quantification
ambiguity, but the primary matrix no longer equates different stable
accessions merely because they share a gene name.

\section*{S2. Aggregation and probability-error definitions}

AUROC, ten-bin ECE, Brier score and log-loss are computed within each
source cohort on its stored out-of-fold predictions. Equal-cohort macro
values average those six metrics. Pooled values concatenate all
predictions before computing the metric. For sample and patient CV,
within-cohort metrics combine predictions from different CV models;
for LOGO, each cohort is scored by one model trained on the other five.
The pooled AUROC additionally includes cross-cohort positive--negative
pairs. It is not the mean of within-cohort or CV-fold AUROCs.

The Brier score is $n^{-1}\sum_i(p_i-y_i)^2$. Log-loss is
$-n^{-1}\sum_i[y_i\log p_i+(1-y_i)\log(1-p_i)]$, with numerical clipping
provided by the recorded metric implementation. Ten-bin ECE weights
the absolute mean-probability/tumor-fraction discrepancy by each bin's
sample mass. No causal meaning is assigned to an AUROC increment or to
the AUROC--ECE association across these six confounded cohorts.

\clearpage
\section*{S3. Cohort-resolved primary and comparator results}

The tables below contain one held-out cohort per LOGO fold. ECE uses ten
equal-width probability bins; Brier and log-loss are unbinned complementary
summaries. No historical confidence intervals are attached to these
corrected predictions.

\metrictable{Primary accession-aligned XGBoost abundance predictions.}{primary_rows.tex}
\metrictable{Three reference-normalized global abundance summaries:
observed mean, median and detection fraction. These are not raw total
intensity or an identified biological component.}{global_rows.tex}
\metrictable{Missingness-only XGBoost predictions on the same accession
matrix. Weak within-cohort ranking and pooled ranking are distinct.}{missingness_rows.tex}

\clearpage
\section*{S4. Mapping, model-family and imputation sensitivities}

These analyses use the corrected tissue and patient identities. The
mapping sensitivity changes the feature set; the L2 and deliberately
leaky-median arms retain the primary 8,906 accessions and six LOGO folds.
They are sensitivities on the same source cohorts, not independent
replications.

\metrictable{The 8,441 first-accession-consistent-gene mapping sensitivity.
No inconsistent historical gene/protein match is retained.}{mapping_rows.tex}
\metrictable{L2 logistic regression with train-fold median imputation and
train-fold standardization.}{l2_rows.tex}
\metrictable{XGBoost with deliberately full-dataset medians, including
held-out feature values but no held-out labels. This is a diagnostic,
not the recommended evaluation.}{leaky_rows.tex}

\clearpage
\section*{S5. Conditional paired-patient resampling}

True patients are sampled with replacement within each of the six
fixed cohorts, carrying all of each patient's sample predictions together.
The same patient draw is used for sample CV, patient CV and LOGO, so
each contrast is paired. We use 2,000 draws and NumPy random seed
20260922. The stored model predictions are not refitted.

The 2.5th--97.5th percentile ranges in Table~\ref{tab:resampling}
describe sensitivity to the observed patient mixture conditional on
those predictions and cohorts. They exclude retraining uncertainty,
model-selection variability and variation in future study environments.
The source cohorts are not asserted to be exchangeable random studies.
Thus these ranges are not full algorithm-refit or future-study confidence
intervals.

\begin{table}[htbp]\centering
\caption{Paired protocol AUROC contrasts and conditional percentile ranges.}
\label{tab:resampling}
\begin{tabular}{llrr}\toprule
Aggregation & Contrast & Point & 2.5th--97.5th percentiles \\
\midrule\tableinput{tables/resampling_rows.tex}\bottomrule
\end{tabular}\end{table}

\section*{S6. Linear model and recalibration implementation}

The L2 classifier is \texttt{LogisticRegression} with $C=1$,
\texttt{solver=lbfgs}, L2 penalty and \texttt{max\_iter=500}.
Missing values use training-fold medians and scaling parameters are
estimated on the training fold only. All six fits finished without a
convergence warning (Table~\ref{tab:convergence}). The superseded
method-lock text giving 1000 iterations does not describe these fits.

\begin{table}[htbp]\centering
\caption{Recorded convergence of the six L2 LOGO fits.}
\label{tab:convergence}
\begin{tabular}{lrrl}\toprule
Held-out cohort & Iterations & Limit & Converged \\
\midrule\tableinput{tables/l2_convergence_rows.tex}\bottomrule
\end{tabular}\end{table}

For recalibration, the recorded cohort order is CCRCC, LUAD, UCEC,
HNSCC, LSCC and PDAC. After removing the outer test cohort, the first
two remaining cohorts and the other three form the inner groups.
Training on one and scoring the other, then swapping, yields the
calibration data. Inner training-fold imputation is separate from the
outer fit. Temperature scaling uses probability logits clipped to
$[10^{-7},1-10^{-7}]$ before transformation, and bounded NLL
minimization over $T\in[0.05,20]$. Platt scaling uses logistic
regression on those logits ($C=10^6$, maximum 2000 iterations).
The maps never use outer held-out labels.

The inner models train on two or three cohorts, unlike the outer
five-cohort predictor. This training-size mismatch limits generalization
of any conclusion about recalibration. Distinct cohort-specific maps
can also change cross-cohort pooled rankings while preserving each
within-cohort monotone ranking. The main article reports the observed
benefits and losses of this specific procedure without a universal claim.

\clearpage
\section*{S7. Side-study dependency disposition}

The following distinctions prevent a historical computation from being
mistaken for evidence on the corrected matrix. Archived files remain
available; removal of their numerical claims from the current argument
does not imply that their scientific questions should be abandoned.

\textbf{Corrected matrix sensitivities retained.}
The 8,441-feature mapping comparison, full-data-median XGBoost,
train-fold L2, coverage curves, and split-half recalibration were rerun
on corrected labels and aligned proteins, and verified against their
stored prediction rows. Their results appear in the article and this
supplement.

\textbf{OV plex-LOGO retained as historical summary context.}
The ovarian table and TMT annotation are independent of the six-cohort
intersection. Their stored source checksums and the 103-sample,
84-tumor/19-normal, 12-plex annotation inventory were checked during
integration. The recorded pooled abundance AUROC is 0.985 under
plex-LOGO and 0.987 under random CV; plex-LOGO ECE is 0.046 and
global-summary AUROC is 0.951. There are only one to three normals per
plex. These are summary-level historical results, not a new
prediction-level replay. They are not promoted as an independently verified
counterexample and do not deconfound the primary cancer/study comparison.
The current six-cohort evidence itself does not support a universal gap.
Historical bootstrap endpoints are
not reused as uncertainty for the corrected six-cohort result.

\textbf{Cohort-identity and label-permutation diagnostics not promoted.}
The old six-class identity classifier, one-versus-rest XGBoost check,
training-cohort identity estimates, and within-/across-cohort label
permutations use the old feature/label definitions. Their numerical
claims require a corrected-object rerun before reuse. They are not
needed to establish the observed primary protocol contrast.

\textbf{Location/scale and representation diagnostics not promoted.}
Per-cohort median/MAD, mean/SD and ComBat-style location/scale restores,
rank, quantile-normal, SVD whitening and PC-removal analyses depend on
the previous matrix. Some additionally use held-out feature moments
and therefore evaluate a target-feature-aware procedure. None of their
historical numbers is presented as a corrected primary estimate.

\textbf{Residual and injected-shift constructions not promoted.}
Within-cohort, within-class mean removal uses outcome labels and is a
diagnostic construction, not a deployable held-out prediction procedure.
Synthetic label-null injection estimates a response to a specified
imposed shift. It is not a biological decomposition, and its previous
target magnitude depended on an obsolete headline AUROC. Both would
need redefinition and rerunning to support new numerical claims.

\textbf{Two-way variance claim withdrawn.}
The earlier biological/cohort/interaction allocation does not identify
biological and acquisition fractions in a confounded design. Its
pseudo-interaction calculation is not retained as an interpretable
two-way variance decomposition; rerunning the same formula on new
data would not cure that problem.

\textbf{BCM/UMICH dual-platform numbers not promoted.}
The three-cohort comparison uses a distinct gene-level mapping and
same-patient join. It is neither an independent-patient validation nor
a crossed cancer-by-batch design. A separate accession/aggregation and
sample-identity audit is required before quantitative reuse of its
effect-size correlations and LOGO estimates. No corrected six-cohort
claim relies on those archived numbers.

\textbf{Other acquisition-metadata screens.}
BRCA plex labels, the unavailable COAD mapping, and GBM/other cohort
eligibility screens are historical source-inventory observations,
not new predictive validation. No numerical BRCA or COAD performance
claim is made. No pathway-enrichment analysis was performed.

\section*{S8. Reproducibility boundary}

The corrected fits use \texttt{cptac} 1.5.14, XGBoost 3.3.0, NumPy
2.5.0, pandas 3.0.3 and scikit-learn 1.9.0. The primary model keeps the
historical hyperparameter settings without an explicitly set XGBoost
random state; sample stratification uses seed 28. Current figures and
TeX tables derive from the verified stored predictions, so rendering
does not invoke a model fit. Hash checks establish version identity,
not scientific validity.

The old submission directories are historical, unsynchronized packages.
They must not be uploaded as this revised article. Full scientific
clearance, any reinstatement of dependent side-study claims, and
submission-package regeneration are separate from a successful local
PDF build.
\end{document}
